\subsection{BÀI TẬP TRẮC NGHIỆM}
\ind{PHẦN I.} \inden{Câu trắc nghiệm nhiều phương án lựa chọn. Mỗi câu hỏi học sinh chỉ chọn một phương án.}\\
\setcounter{ex}{0}
\Opensolutionfile{ans}[ans/1D7-B2-1]

\begin{ex}%[1D5Y2-1]%
	Cho $ u=u(x), v=v(x), v(x)\ne 0 $. Hãy chọn khẳng định \textbf{sai}?
	\choice
	{$(u+v)'=u'+v'$}
	{$(ku)'=ku'$}
	{\True $\left(\dfrac{1}{v}\right)'=-\dfrac{v'}{v}$}
	{$(uv)'=u'v+uv'$}
	\loigiai{
		Ta có  $\left(\dfrac{1}{v}\right)'=-\dfrac{v'}{v^2}$.
	}
\end{ex}

\begin{ex}
	Đạo hàm của hàm số $y=5x^3-x^2-1$ là
	\choice
	{\True $15x^2-2x$}
	{$15x^2-2x$}
	{$15x^2-2x-1$}
	{$-2x$}
	\loigiai{
		TXĐ: $\mathscr{D}=\mathbb{R}$.\\
		Ta có: $y'=\left(5x^3-x^2-1 \right)'=15x^2-2x$.
	}
\end{ex}

\begin{ex}%[1D5Y3-1]%
	Tính đạo hàm của hàm số $y=\sin 2x$ .
	\choice
	{$y'=-\cos 2x$}
	{$y'=-2\cos 2x$}
	{$y'=\cos 2x$}
	{\True $y'=2\cos 2x$}
	\loigiai{
		Ta có $y'=(2x)'\cdot \cos 2x=2\cos 2x$.}
\end{ex}


\begin{ex}
	Tính đạo hàm của hàm số $ y=2018^x $.
	\choice
	{$ y'=x\cdot 2018^{x-1} $}
	{$ y'=\dfrac{2018^x}{\ln 2018} $}
	{\True $ y'=2018^x \cdot \ln 2018 $}
	{$ y'=2018^x $}
	\loigiai{
		Ta có  $ y'=(2018^x)'=2018^x \cdot \ln 2018$.
	}
\end{ex}

\begin{ex}
	Cho hàm số $y=\ln x$. Tính đạo hàm của hàm số trên khoảng $(0; +\infty)$.
	\choice
	{$y'=\dfrac{1}{x\ln10}$}
	{$y'=x$}
	{\True $y'=\dfrac{1}{x}$}
	{$y'=-\dfrac{1}{x}$}
	\loigiai{
		Ta có $y'=(\ln x)'=\dfrac{1}{x}$.
	}
\end{ex}

\begin{ex}
	Đạo hàm của hàm số $f(x)=\log_2 x$ là
	\choice
	{$x\ln2$}
	{$\dfrac{\ln2}{x}$}
	{\True $\dfrac{1}{x\ln2}$}
	{$\dfrac{x}{\ln2}$}
	\loigiai{
		Đạo hàm của hàm số $f(x)=\log_2 x$ là $f'(x)=\dfrac{1}{x\ln2}$.
	}
\end{ex}

\begin{ex}
	Tìm đạo hàm của hàm số $y=\dfrac{2x-1}{x+2}$
	\choice
	{$y'=\dfrac{-5}{(x+2)^2}$}
	{$y'=\dfrac{2}{(x+2)^2}$}
	{$y'=\dfrac{3}{(x+2)^2}$}
	{\True $y'=\dfrac{5}{(x+2)^2}$}
	\loigiai
	{
		Ta có $y'=\dfrac{(2x-1)'(x+2)-(x+2)'(2x-1)}{(x+2)^2}=\dfrac{2x+4-2x+1}{(x+2)^2}=\dfrac{5}{(x+2)^2}.$
	}
\end{ex}

\begin{ex}%[1D5B2-1]
	Tính đạo hàm của hàm số $y=\dfrac{x^2+2x-3}{x+2}.$
	\choice
	{\True  $y'=1+\dfrac{3}{(x+2)^2}$}
	{$y'=\dfrac{x^2+6x+7}{(x+2)^2}$}
	{$y'=\dfrac{x^2+4x+5}{(x+2)^2}$}
	{$y'=\dfrac{x^2+8x+1}{(x+2)^2}$} 
	\loigiai{
		Ta có $y=x-\dfrac{3}{x+2}\Rightarrow y'=1+\dfrac{3}{(x+2)^2}$. 
	}
\end{ex}

\begin{ex}
	Đạo hàm của hàm số $y = \log_3 (4x + 1)$ là
	\choice
	{\True $y' = \dfrac{4}{(4x + 1)\ln 3}$}
	{$y' = \dfrac{1}{(4x + 1) \ln 3}$}
	{$y' = \dfrac{\ln 3}{4x + 1}$}
	{$y' = \dfrac{4 \ln 3}{4x + 1}$}
	\loigiai{
		Ta có $y' = \dfrac{(4x + 1)'}{(4x + 1) \ln 3} = \dfrac{4}{(4x + 1) \ln 3}$.}
\end{ex}

\begin{ex}
	Tính đạo hàm của hàm số $f(x)=\mathrm{e}^{2x-3}$.
	\choice
	{$f'(x)=2\cdot\mathrm{e}^{x-3}$}
	{$f'(x)=\mathrm{e}^{2x-3}$}
	{$f'(x)=-2\cdot\mathrm{e}^{2x-3}$}
	{\True $f'(x)=2\cdot\mathrm{e}^{2x-3}$}
	\loigiai{
		Ta có $f'(x)=(2x-3)'\cdot\mathrm{e}^{2x-3}=2\cdot\mathrm{e}^{2x-3}$.
	}
\end{ex}

\begin{ex}%[1D5B2-1]%
	Đạo hàm của hàm số $y=\sqrt{x^2-5x}$ là
	\choice
	{\True $y'=\dfrac{2x-5}{2\sqrt{x^2-5x}}$}
	{$y'=\dfrac{1}{2\sqrt{x^2-5x}}$}
	{$y'=\dfrac{2x-5}{\sqrt{x^2-5x}}$}
	{$y'=-\dfrac{2x-5}{\sqrt{x^2-5x}}$}
	\loigiai
	{
		$$y'=\dfrac{2x-5}{2\sqrt{x^2-5x}}. $$
	}
\end{ex}

\begin{ex}
	Đạo hàm của hàm số $y=\left(x^3-2x^2\right)^2$ bằng
	\choice
	{$6x^5+16x^3$}
	{$6x^5-20x^4+4x^3$}
	{\True $6x^5-20x^4+16x^3$}
	{$6x^5-20x^4-16x^3$}
	\loigiai
	{
		Ta có $y'=2\left(x^3-2x^2\right)\left(x^3-2x^2 \right)'=\left(x^3-2x^2\right)\left(6x^2-8x\right)=6x^5-20x^4+16x^3$.
	}
\end{ex}

\begin{ex}%[TT Chuyên Hùng Vương đợt 2 - Bình Dương]%[2D2K4]
	Cho hàm số $f(x)=\ln \left (3x-x^2\right )$. Tìm tập nghiệm $S$ của phương trình $f'(x) =0$.		
	\choice
	{$S= \varnothing$}
	{\True $S =\left\{ \dfrac{3}{2}\right \}$}
	{$S =\left\{0;3\right \}$}
	{$S =\left( -\infty;0\right ) \cup \left (3;+\infty\right )$}
	\loigiai
	{
		\begin{enumerate}
			\item Tập xác định $\mathbb{D} =(-\infty;0)\cup (3;+\infty)$
			\item Ta có $f'(x) =\dfrac{2x-3}{x^2 -3x}; f'(x) =0\Leftrightarrow x=\dfrac{3}{2}$.
		\end{enumerate}
		Vậy tập nghiệm $S =\left\{ \dfrac{3}{2}\right \}$.
	}
\end{ex}

\begin{ex}%[1D7H2-1]
	Cho hàm số $y=x^3-3x+2017$. Bất phương trình $y' < 0$ có tập nghiệm là
	\choice
	{\True $S=\left(-1;1\right)$}
	{$S=\left(-\infty;-1\right)\cup \left(1;+\infty \right)$}
	{$\left(1;+\infty \right)$}
	{$\left(-\infty;-1\right)$}
	\loigiai{
		\begin{itemize}
			\item $y=x^3-3x+2017\Rightarrow y'=3x^2-3$;
			\item $y' < 0\Leftrightarrow x^2-1 < 0\Leftrightarrow-1 < x < 1$.
		\end{itemize}
	}
\end{ex}

\begin{ex}
	Cho hàm số $ f(x)=\sqrt{-5x^2+14x-9} $. Tập hợp các giá trị của $ x $ để $ f'(x)<0 $ là
	\choice
	{\True $ \left(\dfrac{7}{5};\dfrac{9}{5}\right) $}
	{$ \left(\dfrac{7}{5};+\infty\right) $}
	{$ \left(-\infty;\dfrac{7}{5}\right) $}
	{$ \left(1;\dfrac{7}{5}\right) $}
	\loigiai{
		Tập xác định: $ \mathscr{D}=\left[1;\dfrac{9}{5}\right] $.
		\\
		Ta có $ f'(x)=\dfrac{-5x+7}{\sqrt{-5x^2+14x-9}} $.
		\\
		$ f'(x)<0\Leftrightarrow \heva{&-5x+7<0\\&-5x^2+14x-9>0 }\Leftrightarrow x\in \left(\dfrac{7}{5};\dfrac{9}{5}\right) $.
	}
\end{ex}

\begin{ex}%Câu 5.%[1D7H2-3]
	Viết phương trình tiếp tuyến của đồ thị hàm số $y=x^4-4x^2+5$ tại điểm có hoành độ $x=-1$. 
	\choice
	{$y=4x-6$}
	{$y=4x+2$}
	{\True $y=4x+6$}
	{$y=4x-2$}
	\loigiai{
		Ta có $y'=4x^3-8x$, $y'(-1)=4$.\\
		Điểm thuộc đồ thị đã cho có hoành độ $x=-1$ là $M(-1;2)$.\\
		Vậy phương trình tiếp tuyến của đồ thị hàm số tại $M(-1;2)$ là
		$$y=y'(-1)(x+1)+2\Leftrightarrow y=4(x+1)+2\Leftrightarrow y=4x+6.$$}
\end{ex}

\Closesolutionfile{ans}

\ind{PHẦN II.} \inden{Câu trắc nghiệm đúng sai. Trong mỗi ý a), b), c), d) ở mỗi câu, học sinh chọn đúng hoặc sai.}\\
\Opensolutionfile{ans}[ans/1D7-B2-2]

\begin{ex}%Câu 1.%[1D7V2-4]
	Cho hàm số $y=f(x)=\dfrac{3x+2}{x+1}$. Các mệnh đề sau đúng hay sai?
	\choiceTF
	{\True $f'(x)>0$, $\forall x\neq-1$}
	{$f'(2)+f'(-3)=\dfrac{5}{36}$}
	{Tích tất cả các nghiệm của phương trình $f'(x)=25$ là $\dfrac{26}{25}$}
	{\True Tiếp tuyến của đồ thị hàm số tại điểm $M(-2;4)$ đi qua điểm $E(-5;1)$}
	\loigiai{
		\begin{itemchoice}
			\itemch \textbf{Đúng}.\\
			$f'(x)=\dfrac{(3x+2)'\cdot (x+1)-(3x+2)\cdot (x+1)'}{(x+1)^2}=\dfrac{1}{(x+1)^2}>0$, $\forall x\neq-1$.
			\itemch \textbf{Sai}.\\
			Ta có $f'(2)+f'(-3)=\dfrac{1}{(2+1)^2}+\dfrac{1}{(-3+1)^2}=\dfrac{13}{36}$.
			\itemch \textbf{Sai}.\\
			Ta có $f'(x)=25\Leftrightarrow\dfrac{1}{(x+1)^2}=25\Leftrightarrow(x+1)^2=\dfrac{1}{25}\Leftrightarrow\hoac{&x_1=-\dfrac{6}{5}\\&x_2=-\dfrac{4}{5}.}$ \\
			Nên $x_1\cdot x_2=\dfrac{24}{25}$.
			\itemch  \textbf{Đúng}.\\
			Ta có $f'(-2)=\dfrac{1}{(-2+1)^2}=1$.\\
			Phương trình tiếp tuyến của đồ thị hàm số tại điểm $M(-2;4)$ là $$y-4=1(x+2)\Leftrightarrow y=x+6.$$
			Vì $1=-5+6$ suy ra tiếp tuyến đi qua điểm $E(-5;1)$.
	\end{itemchoice}}
\end{ex}

\begin{ex} %[Word to LaTeX 3.2]
	Cho hàm số $f(x)=4\sin x+2x+1$.
	\choiceTF
	{\True $f(0)=1;f\left(-\dfrac{\pi}{2}\right)=-\pi -3$}
	{Đạo hàm của hàm số đã cho là $f'(x)=-4\cos x+2$}
	{\True Nghiệm của phương trình $f'(x)=0$ trên đoạn $\left[{0;\pi}\right]$ là $\dfrac{2\pi}{3}$}
	{\True Hệ số góc tiếp tuyến của đồ thị đã cho có giá trị lớn nhất bằng $6$.}
	\loigiai{
		\begin{itemchoice}
			\itemch Đúng.\\ Vì $f(0)=1;f\left(-\dfrac{\pi}{2}\right)=-\pi -3$
			\itemch Sai.\\ Vì $f'(x)=4\cos x+2$
			\itemch Đúng. \\Ta có $f'(x)=4\cos x+2$.\\
			Xét $f'(x)=0$$\Leftrightarrow 4\cos x+2=0\Leftrightarrow \cos x=-\dfrac{1}{2}\Leftrightarrow x=\dfrac{2\pi}{3}$, do $x\in \left[{0;\pi}\right]$.
			\itemch Sai.\\ Xét hàm số $f(x)$ trên $\left[{0;\pi}\right]$.\\
			Ta có $f'(x)=4cosx+2, f'(x)=0$ có nghiệm trên $\left[{0;\pi}\right]$ là $x=\dfrac{2\pi}{3}$.\\
			Ta có Ta cso $f'(x)=4\cos x +2 \le 6$. Suy ra hệ số góc của tiếp tuyến có giá trị lớn nhất bằng $6$.
		\end{itemchoice}
	}
\end{ex}

\begin{ex}%[2D2B4-4]%
	Cho hàm số $f(x)=x^2-\ln(1-2x)$. Xét tính đúng sai của các khẳng định sau:
	\choiceTF
	{Tập xác định của hàm số là $(-\infty;2)$}
	{\True Đạo hàm của hàm số là $f'(x)=2x+\dfrac{2}{1-2x}$}
	{Phương trình $f'(x)=0$ có hai nghiệm phân biệt}
	{\True Tiếp tuyến của đồ thị hàm số đã cho tại gốc tọa độ có hệ số góc bằng $2$}
	\loigiai{
		\begin{itemchoice}
			\itemch Hàm số xác định khi $1-2x>0 \Leftrightarrow x<\dfrac{1}{2}$.
			\itemch Ta có $f'(x)=2x+\dfrac{2}{1-2x}$
			\itemch $f'(x)=0\Rightarrow2x+\dfrac{2}{1-2x}=0\Leftrightarrow\hoac{&x=1 \text{ loại }\\&x=-\dfrac{1}{2}}$
			\itemch Ta có $f'(x)=2x+\dfrac{2}{1-2x}$, suy ra $f(0)=2$.
		\end{itemchoice}
	}
\end{ex}

\begin{ex}
	Một vật chuyển động thẳng không đều xác định bởi phương trình $s(t)=t^{2} - 6 t + 1$, trong đó ${s}$ tính bằng mét và ${t}$ tính bằng giây. Xét tính đúng sai của các khẳng định sau
	\choiceTF
	{ Gia tốc chuyển động của vật tại thời điểm $t=10$ là ${6}$ m/$s^2$ }
	{ \True Vận tốc chuyển động của vật tại thời điểm $t=7$ là ${8}$ m/s }
	{ \True Quãng đường vật đi được sau ${7}$ giây kể từ khi bắt đầu chuyển động là ${8}$  m. }
	{ Vận tốc nhỏ nhất vật đạt được trong khoảng thời gian từ $t=3$ đến $t=5$ là $4$ m/s }
	\loigiai{ 
		a) Khẳng định đã cho là sai.
		
		$v(t)=\left(t^{2} - 6 t + 1\right)'=2 t - 6.$
		
		$a(t)=\left(2 t - 6\right)'=2$.
		
		$a(10)=2$ m/$s^2$.\\ 
		b) Khẳng định đã cho là khẳng định đúng.
		
		$v(t)=\left(t^{2} - 6 t + 1\right)'=2 t - 6.$
		
		$v(7)=8$ m/s.\\ 
		c) Khẳng định đã cho là đúng.
		
		$s(7)=8$ m.\\ 
		d) Khẳng định đã cho là khẳng định sai.
		
		$v(t)=\left(t^{2} - 6 t + 1\right)'=2 t - 6.$ Hàm $v(t)$ là hàm số đồng biến trên $\mathbb{R}$.
		
		Trong khoảng thời gian từ $t=3$ đến $t=5$ thì vận tốc đạt nhỏ nhất tại $t=3$.
		
		Vận tốc đạt được khi đó là $v\left(3\right)=0$.
		
}\end{ex}


\Closesolutionfile{ans}

\ind{PHẦN III.} \inden{Câu trắc nghiệm trả lời ngắn.}\\
\Opensolutionfile{ans}[ans/1D7-B2-3]
\begin{ex}%[Đỗ Minh Phúc]%[1D5Y2-6]%
	Một vật chuyển động theo quy luật $s=-\dfrac{1}{2}t^2+20t$ với $t$ (giây) là khoảng thời gian tính từ khi vật bắt đầu chuyển động và $s$ (mét) là quãng đường vật đi được trong thời gian đó. Hỏi vận tốc tức thời của vật tại thời điểm $t=8$ giây bằng bao nhiêu m/s?
	\\
	\shortans[oly]{$12$}
	\loigiai{
		Vận tốc tức thời được tính bằng công thức $v(8)= \left(s(8)\right)^{'}= 12$ m/s.}
\end{ex}


\begin{ex}%[1D5B2-6]
	Nếu số lượng sản phẩm sản xuất được của một nhà máy là $x$ (đơn vị: trăm sản phẩm) thì lợi nhuận sinh ra là $P(x)=-200x^2+ 12800x-74000$ (nghìn đồng). Tính tốc độ thay đổi lợi nhuận của nhà máy đó khi sản xuất  $2000$ sản phẩm.
	\\
	\shortans[oly]{$4800$}
	\loigiai{
		Ta có $P'(x)=-2 \cdot 200 x + 12800 = -400x+12800$.\\
		Tốc độ thay đổi lợi nhuận của nhà máy đó khi sản xuất $1200$ sản phẩm là $P'(20)= -400 \cdot 20 + 12800 = 4800$.
	}
\end{ex}

\begin{ex}%[1D7C2-8]
	Một con lắc lò xo dao động điều hoà theo phương ngang trên mặt phẳng không ma sát, có phương trình chuyển động $x=4\cos \left(\pi t-\dfrac{2\pi}{3} \right)+4$ (cm), trong đó $t$ là thời gian tính bằng giây. Tìm thời điểm mà vận tốc tức thời của con lắc bằng $0$ (cm/s) là $t=\dfrac{a}{b}+k(k\in \mathbb{Z})$ (s), trong đó $a,b$ là các số nguyên và phân số $\dfrac{a}{b}$ là phân số tối giản. Tính tổng $a+b$.\\
	\shortans[oly]{$5$}
	\loigiai{
		Ta có $v(t)=x'=-4\pi \sin \left(\pi t-\dfrac{2\pi}{3} \right)$.\\	
		Thời điểm mà vận tốc tức thời của con lắc bằng $0$ nghĩa là $v(t)=0$.
		\begin{eqnarray*}
			&\Leftrightarrow& -4\pi \sin \left(\pi t-\dfrac{2\pi}{3} \right)=0\\
			&\Leftrightarrow& \sin \left(\pi t-\dfrac{2\pi}{3} \right)=0\\
			&\Leftrightarrow& \pi t-\dfrac{2\pi}{3}=k\pi\\
			&\Leftrightarrow& t=\dfrac{2}{3}+k,\,\left(k\in\mathbb{Z}\right).
		\end{eqnarray*}
		Vậy các thời điểm mà vận tốc tức thời của con lắc bằng $0$ là $t=\dfrac{2}{3}+k(k\in \mathbb{Z})(s)$.\\
		Suy ra $a=2, b=3$, nên $a+b=5$.
	}
\end{ex}

\begin{ex}%[1D7C2-3]
	Cho hàm số $y=\dfrac{9}{x}$ có đồ thị là $(C)$. Biết tiếp tuyến của đồ thị $(C)$ tại điểm $M(3;3)$ tạo với hai trục toạ độ một tam giác. Tính diện tích tam giác đó.\\
	\shortans[oly]{$18$}
	\loigiai{
		\immini{Ta có $y'=-9\cdot \dfrac{1}{x^2}$.\\
			Hệ số góc của tiếp tuyến tại điểm $M$ là $y'(3)=\dfrac{-9}{3^2}=-1$.\\
			Phương trình tiếp tuyến $(\Delta)$ với $(C)$ tại tiếp điểm $M$ là
			\[y-3=-1(x-3)\Leftrightarrow y=-x+6.
			\]
			Biết $(\Delta)$ cắt trục hoành và trục tung lần lượt tại hai điểm $A(6;0),B(0;6)$ nên diện tích tam giác $OAB$ vuông tại $O$ bằng
			$S_{\triangle OAB}=\dfrac{1}{2} OA\cdot OB=\dfrac{1}{2} \cdot 6\cdot 6=18\; \;$ (đơn vị diện tích).}
		{\begin{tikzpicture}[line cap=butt,line join=miter,>=stealth,scale=.4,font=\footnotesize]
				\tikzset{declare function={xmin=-1;xmax=7;ymin=-1;ymax=7;
						f(\x)=-1*(\x)+6;
					},
					smooth,samples=450
				}
				\draw[->] (xmin,0)--(xmax,0) node[shift={(-100:7pt)}]{$ x $};
				\draw[->] (0,ymin)--(0,ymax) node[shift={(190:7pt)}]{$ y $};
				\draw[fill=white] (0,0) node[shift={(225:9pt)}]{$ O $} circle (1.5pt)
				(6,0) node[shift={(45:7pt)}]{$6$} circle (1.5pt)
				(0,6) node[shift={(45:7pt)}]{$6$} circle (1.5pt);
				\clip (xmin,ymin) rectangle (xmax,ymax);
				\draw  plot[domain=xmin:xmax] (\x, {f(\x)});
		\end{tikzpicture}}
	}
\end{ex}

\begin{ex}
	Cho hàm số $f(x)=\ln \dfrac{2018x}{x+1}$. Biết tổng $S=f'(1)+f'(2) +\cdots +f'(2018)=\dfrac{2018}{n}$. Tìm $n$.\\
	\shortans[oly]{$2019$}
	\loigiai
	{ Ta có $f'(x)= \dfrac{1}{x} -\dfrac{1}{x+1} \quad \forall x >0$.\\
		Suy ra 
		\begin{eqnarray*}
			S&=&f'(1)+f'(2) +\cdots +f'(2018)\\
			&= & \left( 1 -\dfrac{1}{2}\right)+ \left( \dfrac{1}{2}- \dfrac{1}{3}\right)+ \cdots + \left( \dfrac{1}{2018}  -\dfrac{1}{2019}\right)\\
			&= &  1 -\dfrac{2018}{2019}\\
			& = &\dfrac{2018}{2019}.
		\end{eqnarray*}	
	}
\end{ex}


\Closesolutionfile{ans}

